\section{1X}

\subsection{1m}

\bin{\currfiledir/seq/1m}

\subsection{1M}

We play Kaplan-interchange after 1H, i.e. \texttt{1N = 5+S, 1S = 4-S}.
However, 6+S inv will bid 1S initially, then rebid 2S (3S) regardless of opener's rebid.
\bin{\currfiledir/seq/1M_}

\subsubsection{PH}

\bin{\currfiledir/seq/1M_PH}

\subsection{Gazzilli}

Before introducing opener's rebids, let's introduce (our) Gazzilli, which uses 
\texttt{1C - 1DH - 1S}, \texttt{1D - 1HS - 1N}, \texttt{1H - 1SN - 2C}, \texttt{1S - 1N - 2C} as a relay. showing either

\begin{itemize}
    \setlength\itemsep{0pt}
    \item a standard rebid of {1C, 1D, 1H, or 1S}.
    \item 16+
\end{itemize}

Responding to it is simple: \texttt{+1 = 8+}, {\color{blue} except:
\begin{itemize}
    \setlength\itemsep{0pt}
    \item limit raise: \texttt{1M - 1SN; 2C - 3M}
    \item 8-10 5S4H: \texttt{1m - 1S; Gaz - 2H}
    \item oM invite: \texttt{1H - 1S; 2C - 2S} and \texttt{1S - 1N; 2C - 3H}
\end{itemize}}
otherwise, bid is weak and s/o (2N = pick a minor); and {\color{blue} jump is nat ST}. We'll revisit the opener's rebid later.

\subsection{1m - 1M}

\bin{\currfiledir/seq/1m-1M-simple}
\bin{\currfiledir/seq/1m-1M-fit}

The following design may seems asymmetric, but with a reason: if resp already choose to s/o, there's no need for min unbal to even bid.
In this case, the more severe problem is 15-17 6m3M, one of natrual's nightmare.
However, it's not the case when resp is 8+ and unlimited (since 15-17 6m is already GF), where you might want to show your 3-fit for min bal.
\bin{\currfiledir/seq/1m-1M-Gaz}

\subsection{1m - 1N \& 1C - 1S}
\bin{\currfiledir/seq/1m-1SN}

\subsection{1S - 1N \& 1H - 1SN}

\bin{\currfiledir/seq/1M-1N}

\section{1X subseq}

\subsection{2/1, and nebulous 2C}

We use the "2/1 always GF \& nebulous 2C" approach, where all balanced GF hands without 5D/H starts with 2C so that 2D/H promises 5+ card. 
(4441 starts with 2D, since it is tolarable to bid D with a fit anyways)
We use slightly different systems after \texttt{1M - 2C} and \texttt{1M - 2DH}:

\bin{\currfiledir/seq/1M-2X}
\bin{\currfiledir/seq/1M-2DH}
\bin{\currfiledir/seq/1M-2C}

For subseq, the principles are:
\bin{\currfiledir/seq/1M-2X-subseq}

\subsection{subseq after Gazzilli accepted}

After Gaz accepted and opener is 16+, GF has established. Therefore, the main problems are: 

\begin{itemize}
    \setlength\itemsep{0pt}
    \item finding the correct contract: \texttt{4M} or \texttt{3N}, or occationally \texttt{5m}
    \item investigating slam
\end{itemize}

Doesn't this look like 2-over-1? Therefore, I proposed a simple method: 
exchange the lowest new suit with \texttt{2N}, showing a waiting hand, and after that, the responder can freely bids natrually at the two level or \texttt{2N} as a waiting bid.
Similar to 2/1, direct 3X by opener or resp shows extra shape and strength. (caveat: if \texttt{2N} points to some major suit, opener usually shows that. However, if it points to diamonds, I think it promises extra shape or strength.)

\bin{\currfiledir/seq/Gaz_waiting}

\subsection{after Gaz's weak 2m}

Although we say "same as natrual", but we shall clarify a bit:

\bin{\currfiledir/seq/note_1m}

\subsection{2-way}

\bin{\currfiledir/seq/note_2way}

\subsection{1m - 1X - 2N}

\bin{\currfiledir/seq/note_2N}

\subsection{Jacoby 2N}

see {\color{blue} conventions/ST}

\subsection{minor invite}

We have many minor invites. You expect partner to accept with:
not-too-bad min with honor and support, or max (stronger than normal that would accept inv) otherwise. 
For inviter, it is usually lighter than normal invite, even lighter with longer minor.

% \subsubsection{difficult sequences}

% Here I list some sequences that is easily confused (as inv): 

% \bin{\currfiledir/seq/note_confuse}
